\chapter{Structure Determination: \texorpdfstring{$^{13}$C}{13C} NMR, MS and IR Together}\label{ch:b2:structure-determination}

A fragrance laboratory receives a vial of a colourless liquid that smells of jasmine. Nobody has
written a label. A \omterm{def:b2:mass-spec-atomic:spectrum}{mass spectrum} gives its molecular formula, an infrared spectrum its functional
groups, a carbon-13 spectrum its skeleton and a proton spectrum the way the pieces are joined; within an
hour the liquid has a name. This chapter adds \omterm{def:b2:structure-determination:c13}{carbon-13 NMR} to the proton NMR and infrared spectroscopy
of the Year 1 volume and to the \omterm{def:b2:mass-spec-atomic:spectrum}{mass spectrometry} of \cref{ch:b2:mass-spec-atomic}, and turns the four
into one method.

\begin{recall}
The Year 1 volume: chemical shift and shielding, equivalent protons, integration, spin--spin coupling
and the $n + 1$ rule, infrared wavenumbers of functional groups and the fingerprint region, the degree
of unsaturation, solving a structure from three spectra. \Cref{ch:b2:mass-spec-atomic}: the \omterm{def:b2:mass-spec-atomic:molecular-ion}{molecular
ion}, \omterm{def:b2:mass-spec-atomic:isotope-pattern}{isotope patterns}, fragmentation.
\end{recall}

\section{Carbon-13 NMR}

Carbon-12, the main isotope, has no nuclear spin and gives no NMR signal. Carbon-13 has a spin of
$\tfrac12$, like the proton, but it makes up only 1.1~\% of natural carbon and its signal is weaker,
nucleus for nucleus, than that of a proton; a carbon spectrum therefore needs more sample or more scans.
% ledger: iso:C13

\begin{definition}[Carbon-13 NMR]\label{def:b2:structure-determination:c13}
\emph{Carbon-13 NMR}\index{carbon-13 NMR} records the resonances of the \ce{^{13}C} nuclei of a sample.
It is usually run with \emph{broadband decoupling}\index{broadband decoupling}: the protons are
irradiated over their whole range of frequencies during the measurement, which removes every
carbon--proton coupling, so that each kind of carbon gives one single line. \emph{Equivalent
carbons}\index{equivalent carbons}, exchanged by a symmetry of the molecule (or by a fast rotation),
give the same line.
\end{definition}

\begin{proposition}[No carbon--carbon coupling]\label{prop:b2:structure-determination:no-cc-coupling}
Couplings between neighbouring carbon atoms do not appear in a natural-abundance \ce{^{13}C}
spectrum.
\end{proposition}

\begin{proof}
A coupling between two carbons is seen only in molecules where both are \ce{^{13}C}. For two given
neighbouring positions the probability is $a_{13}^2 = 0.011^2 \approx 1.2 \times 10^{-4}$: about one
molecule in eight thousand. The signal of each carbon comes almost entirely from molecules in which its
neighbours are \ce{^{12}C}, invisible to the spectrometer; the coupled satellites are a hundred times
weaker than the noise of an ordinary spectrum.
\end{proof}

\begin{proposition}[Number of lines]\label{prop:b2:structure-determination:signals}
A decoupled \ce{^{13}C} spectrum shows one line per set of \omterm{def:b2:structure-determination:c13}{equivalent carbons}, unless two sets happen to
have the same shift.
\end{proposition}

\begin{proof}[Argument]
\omterm{def:b2:structure-determination:c13}{Equivalent carbons} are in identical surroundings, so they have the same shift; non-equivalent carbons
generally differ. With couplings to protons removed and couplings between carbons absent
(\cref{prop:b2:structure-determination:no-cc-coupling}), nothing splits a line. Accidental overlap
happens when two different carbons have shifts closer than the resolution, which is common for the
carbons of a benzene ring near \qty{128}{ppm}.
\end{proof}

Shifts spread over about \qty{220}{ppm}, twenty times the range of protons, so overlaps are rarer than
in proton spectra, and the shift alone tells the kind of carbon. The chart below shows the shifts
measured for the compounds of this chapter.

\begin{omfigure}
\begin{tikzpicture}
\begin{axis}[width=0.86\linewidth, height=4.6cm, xmin=0, xmax=220, ymin=-0.6, ymax=4.6, x dir=reverse,
  xlabel={$\delta$ (ppm)}, label style={font=\small}, tick label style={font=\scriptsize},
  ytick={0,1,2,3,4}, yticklabels={alkyl C, C--O, aromatic C, ester C=O, ketone C=O},
  xtick={0,20,40,60,80,100,120,140,160,180,200,220}, grid=major, grid style={black!10}]
\addplot[only marks, mark=|, mark size=5pt, thick, omDef] table[x=shift, y=cls] {figdata/out/bachelor-2-nmr-spectra-d.dat};
\end{axis}
\end{tikzpicture}

{\small Measured \ce{^{13}C} shifts of butan-2-one, pentan-2-one, ethyl benzoate and benzyl ethanoate,
sorted by kind of carbon (CDCl$_3$; PubChem data). Ketone carbonyls lie near 209, ester carbonyls near
167--171, \omterm{def:b2:huckel:aromatic}{aromatic} carbons between 128 and 137, carbons bonded to an ester oxygen near 61--66, and
alkyl carbons below 46.}
% figdata: bachelor-2/nmr-spectra.py part d
% ledger: c13:butanone, c13:pentan-2-one, c13:ethylbenzoate, c13:benzylacetate
\end{omfigure}

\begin{proposition}[Heights are not counts]\label{prop:b2:structure-determination:intensities}
In a routine decoupled \ce{^{13}C} spectrum the heights of the lines are not proportional to the
numbers of carbons they represent; carbons without hydrogens give weak lines.
\end{proposition}

\begin{proof}[Argument]
Routine spectra repeat the measurement faster than the slowest carbons relax back to equilibrium, and
carbons without attached protons relax slowly, so their signal is partly saturated. Decoupling also
enhances the signals of carbons carrying protons, by an amount that depends on their surroundings.
In butan-2-one, the carbonyl line is the weakest of the four although it stands for one carbon like
the others. Integration, so useful for protons, is therefore not used for routine carbon spectra.
\end{proof}

\section{DEPT}

\begin{definition}[DEPT]\label{def:b2:structure-determination:dept}
\emph{DEPT}\index{DEPT} (distortionless enhancement by polarisation transfer) is a set of carbon-13
experiments that give the signals of carbons with attached protons a sign or an absence depending on
their number of hydrogens. A \emph{quaternary carbon}\index{quaternary carbon} carries no hydrogen
(bonded to four other atoms, or a carbonyl or \omterm{def:b2:huckel:aromatic}{aromatic} carbon without H); it gives no DEPT signal.
\end{definition}

\begin{proposition}[Reading DEPT]\label{prop:b2:structure-determination:dept}
In DEPT-135, \ce{CH} and \ce{CH3} carbons point up, \ce{CH2} carbons point down and \omterm{def:b2:structure-determination:dept}{quaternary carbons}
are absent. In DEPT-90, only \ce{CH} carbons appear.
\end{proposition}

\admitted

The experiment transfers magnetisation from the protons to the carbon they are bonded to, through a
sequence of radio-frequency pulses whose last angle selects the response; how the pulses do it is
explained in the Year 3 volume. Comparing the decoupled spectrum with DEPT-135 and DEPT-90 sorts every
line into C, CH, \ce{CH2} or \ce{CH3}.

\begin{omfigure}
\pgfplotsset{dept/.style={width=\linewidth, height=2.9cm, ycomb, x dir=reverse, ymin=-1, ymax=1.15,
  ytick=\empty, axis y line=none, axis x line=middle, x axis line style={-}, tick label style={font=\tiny},
  title style={font=\scriptsize, at={(0.0,0.95)}, anchor=north west}, every axis plot/.append style={thick}}}
\begin{minipage}[t]{0.49\linewidth}\centering
{\small butan-2-one}\par
\begin{tikzpicture}
\begin{axis}[dept, xmin=0, xmax=220, xtick={0,50,100,150,200}, title={decoupled}]
\addplot[omDef] table[x=shift, y=dec] {figdata/out/bachelor-2-nmr-spectra-a.dat};
\end{axis}
\end{tikzpicture}\par
\begin{tikzpicture}
\begin{axis}[dept, xmin=0, xmax=220, xtick={0,50,100,150,200}, title={DEPT-135}]
\addplot[omThm] table[x=shift, y=d135] {figdata/out/bachelor-2-nmr-spectra-a.dat};
\end{axis}
\end{tikzpicture}\par
\begin{tikzpicture}
\begin{axis}[dept, xmin=0, xmax=220, xtick={0,50,100,150,200}, title={DEPT-90}]
\addplot[omMeth] table[x=shift, y=d90] {figdata/out/bachelor-2-nmr-spectra-a.dat};
\end{axis}
\end{tikzpicture}
\end{minipage}\hfill
\begin{minipage}[t]{0.49\linewidth}\centering
{\small ethyl benzoate}\par
\begin{tikzpicture}
\begin{axis}[dept, xmin=0, xmax=220, xtick={0,50,100,150,200}, title={decoupled}]
\addplot[omDef] table[x=shift, y=dec] {figdata/out/bachelor-2-nmr-spectra-b.dat};
\end{axis}
\end{tikzpicture}\par
\begin{tikzpicture}
\begin{axis}[dept, xmin=0, xmax=220, xtick={0,50,100,150,200}, title={DEPT-135}]
\addplot[omThm] table[x=shift, y=d135] {figdata/out/bachelor-2-nmr-spectra-b.dat};
\end{axis}
\end{tikzpicture}\par
\begin{tikzpicture}
\begin{axis}[dept, xmin=0, xmax=220, xtick={0,50,100,150,200}, title={DEPT-90}]
\addplot[omMeth] table[x=shift, y=d90] {figdata/out/bachelor-2-nmr-spectra-b.dat};
\end{axis}
\end{tikzpicture}
\end{minipage}

{\small Decoupled \ce{^{13}C} spectra (measured shifts and heights, PubChem data) with the \omterm{def:b2:structure-determination:dept}{DEPT}
responses (signs from \cref{prop:b2:structure-determination:dept}, heights schematic). Butan-2-one: C=O
209.3 (absent in \omterm{def:b2:structure-determination:dept}{DEPT}), \ce{CH2} 36.9 (down), \ce{CH3} 29.4 and 7.9 (up). Ethyl benzoate: C=O 166.5 and
the ring carbon bearing the ester, 130.6, vanish; three \omterm{def:b2:huckel:aromatic}{aromatic} CH (132.8, 129.6, 128.3) stay in
DEPT-90; \ce{OCH2} 60.9 points down.}
% figdata: bachelor-2/nmr-spectra.py parts a, b
\end{omfigure}

\section{Combining the techniques}

Each technique answers its own question, and the answers are assembled in a fixed order: the formula
first, because it bounds everything else; then the functional groups and the skeleton; then the
connections.

\begin{method}[Solving an unknown]\label{met:b2:structure-determination:solve}
\begin{enumerate}
  \item Formula: \omterm{def:b2:mass-spec-atomic:molecular-ion}{molecular ion}, $M+1$ for the number of carbons, \omterm{def:b2:mass-spec-atomic:isotope-pattern}{isotope patterns}, \omterm{def:b2:mass-spec-atomic:nitrogen-rule}{nitrogen rule}; an
        exact mass if available. Degree of unsaturation.
  \item Infrared: \ce{C=O} (and which kind), \ce{O-H}, \ce{N-H}, \ce{C#N}, \omterm{def:b2:huckel:aromatic}{aromatic} or alkene
        \ce{C-H} above \qty{3000}{cm^{-1}}.
  \item Carbon-13 and \omterm{def:b2:structure-determination:dept}{DEPT}: number of kinds of carbon (symmetry), their classes from the shifts, and
        C, CH, \ce{CH2}, \ce{CH3} from \omterm{def:b2:structure-determination:dept}{DEPT}. Compare with the formula.
  \item Proton NMR: integrations, shifts, multiplicities; neighbours from the $n + 1$ rule.
  \item Assemble the fragments; write every candidate structure.
  \item Check each candidate against every datum, including the mass-spectrum fragments; keep the one
        that explains them all.
\end{enumerate}
\end{method}

\begin{omfigure}
\begin{tikzpicture}[x=1cm, y=1cm, st/.style={draw=black!70, rounded corners=2pt, fill=white,
  font=\small, align=center, minimum height=0.9cm, inner sep=3pt, text width=2.3cm}]
\node[st, fill=omDef!10] (ms) at (0,0) {MS\\formula, $M+1$};
\node[st] (dbe) at (2.9,0) {degree of\\unsaturation};
\node[st, fill=omDef!10] (ir) at (5.8,0) {IR\\functional groups};
\node[st, fill=omDef!10] (c13) at (8.7,0) {\ce{^{13}C}, DEPT\\skeleton};
\node[st, fill=omDef!10] (h1) at (8.7,-1.8) {\ce{^1H}\\connections};
\node[st] (cand) at (5.8,-1.8) {candidate\\structures};
\node[st, fill=omThm!10] (chk) at (2.9,-1.8) {check every\\datum};
\node[st] (ans) at (0,-1.8) {structure};
\draw[->, thick] (ms) -- (dbe); \draw[->, thick] (dbe) -- (ir); \draw[->, thick] (ir) -- (c13);
\draw[->, thick] (c13) -- (h1); \draw[->, thick] (h1) -- (cand); \draw[->, thick] (cand) -- (chk);
\draw[->, thick] (chk) -- (ans);
\draw[->, thick, densely dashed] (chk.south) -- ++(0,-0.5) -| node[pos=0.25, below, font=\footnotesize]
  {a datum unexplained: back to the candidates} (cand.south);
\end{tikzpicture}
\par\vspace{4pt}

{\small The order of work for an unknown (\cref{met:b2:structure-determination:solve}). The last step
is the one most often skipped: a structure that leaves one peak unexplained is not yet the answer.}
\end{omfigure}

\section{Worked cases}

\begin{example}[A ketone, C$_5$H$_{10}$O]\label{ex:b2:structure-determination:ketone}
An unknown has $M = 86$, an IR band near \qty{1715}{cm^{-1}}, carbon lines at 208.9, 45.7, 29.8, 17.4 and
13.7~ppm, and proton signals at 2.41 (t, 2H), 2.13 (s, 3H), 1.61 (sextet, 2H) and 0.91 (t, 3H).
One degree of unsaturation, a saturated ketone (no aldehyde proton near 9.7). Five carbons, five lines:
no symmetry. The singlet of three protons at 2.13 is a \ce{CH3} next to the carbonyl; the triplet,
sextet, triplet sequence is a propyl group whose \ce{CH2} at 2.41 touches the carbonyl. Pentan-2-one,
\ce{CH3COCH2CH2CH3}. Its symmetric isomer pentan-3-one would show three carbon lines and no singlet;
its \omterm{def:b2:mass-spec-atomic:spectrum}{mass spectrum} has the McLafferty ion at 58 that pentan-3-one lacks (\cref{ch:b2:mass-spec-atomic}).
% ledger: c13:pentan-2-one, nmr:pentan-2-one, ir:CO-ketone, ms:pentan-2-one
\end{example}

\begin{example}[An aromatic ester, C$_9$H$_{10}$O$_2$]\label{ex:b2:structure-determination:ester}
An unknown of formula \ce{C9H10O2} (five degrees of unsaturation) shows a carbonyl band and seven
carbon lines: 166.5 (C), 132.8 (CH), 130.6 (C), 129.6 (CH), 128.3 (CH), 60.9 (\ce{CH2}) and 14.3
(\ce{CH3}). Its proton spectrum: 8.05 (m, 2H), 7.35--7.63 (m, 3H), 4.37 (q, 2H), 1.38 (t, 3H). A
monosubstituted benzene ring (four \omterm{def:b2:huckel:aromatic}{aromatic} carbon lines, two of them for two carbons each; five
\omterm{def:b2:huckel:aromatic}{aromatic} protons) accounts for four unsaturations, the carbonyl for the fifth. The quartet--triplet pair
is an ethyl group; its \ce{CH2} at 4.37 (and at 60.9 in carbon) is bonded to oxygen. The ester carbonyl
at 166.5 is bonded to the ring: the two protons at 8.05, next to it, are pushed downfield. Ethyl
benzoate, \ce{C6H5COOCH2CH3}.
% ledger: c13:ethylbenzoate, nmr:ethylbenzoate
\end{example}

\section{Traps}

\begin{itemize}
  \item \emph{Exchangeable protons} (\ce{O-H}, \ce{N-H}) give broad signals at variable shifts, are
        often not coupled to their neighbours, and disappear when a drop of \ce{D2O} is shaken with the
        sample.
  \item \emph{Overlapping signals}: two carbons of a benzene ring, or several \ce{CH2} of a chain, may
        share one line; count lines against the formula before concluding on symmetry.
  \item \emph{Diastereotopic protons}: the two protons of a \ce{CH2} next to a stereocentre are not
        equivalent, may have different shifts and couple to each other.
  \item \emph{Second-order patterns}: when two coupled protons have shifts closer than a few times their
        coupling constant, the multiplets lean towards each other and the $n+1$ rule fails; a
        higher-field spectrometer restores simple patterns.
\end{itemize}

\begin{safety}{\ghs[5mm]{GHS02}\,\ghs[5mm]{GHS07}\,\ghs[5mm]{GHS08}}
Butan-2-one is flammable. The magnet of an NMR spectrometer is always on: no steel tools, gas
cylinders or magnetic cards near it, and no access for people with pacemakers.
% ledger: ghs:butanone, ghs:benzylacetate
\end{safety}

\section{Exercises}

\begin{exercise}[$\star$]\label{exo:b2:structure-determination:1}
How many lines does the decoupled \ce{^{13}C} spectrum of each xylene (1,2-, 1,3- and
1,4-dimethylbenzene) show?
\end{exercise}

\begin{exercise}[$\star$]\label{exo:b2:structure-determination:2}
Predict the DEPT-135 and DEPT-90 spectra of butan-2-ol.
\end{exercise}

\begin{exercise}[$\star$]\label{exo:b2:structure-determination:3}
To which kind of carbon does a line at 209, 170, 129, 64 and 14~ppm most likely belong?
\end{exercise}

\begin{exercise}[$\star$]\label{exo:b2:structure-determination:4}
Compute the degree of unsaturation of \ce{C8H8O2} and propose a structure containing a benzene ring.
\end{exercise}

\begin{exercise}[$\star\star$]\label{exo:b2:structure-determination:5}
A liquid \ce{C4H8O} shows an IR band at \qty{1715}{cm^{-1}} and carbon lines at 209, 37, 29 and 8~ppm
(DEPT-135: 37 down, 29 and 8 up, 209 absent). Identify it.
\end{exercise}

\begin{exercise}[$\star\star$]\label{exo:b2:structure-determination:6}
How would you tell pentan-2-one from pentan-3-one with \ce{^{13}C} NMR alone? With proton NMR? With
\omterm{def:b2:mass-spec-atomic:spectrum}{mass spectrometry}?
\end{exercise}

\begin{exercise}[$\star\star$]\label{exo:b2:structure-determination:7}
An ester has $M = 88$ and proton signals at 4.12 (q, 2H), 2.04 (s, 3H) and 1.26 (t, 3H) (exercise
data). Identify it.
\end{exercise}

\begin{exercise}[$\star\star$]\label{exo:b2:structure-determination:8}
Distinguish 1,2-dichlorobenzene from 1,4-dichlorobenzene by \ce{^{13}C} NMR.
\end{exercise}

\begin{exercise}[$\star\star$]\label{exo:b2:structure-determination:9}
In 2-chlorobutane, why are the two protons of the \ce{CH2} group not equivalent? What does that do to
the proton spectrum?
\end{exercise}

\begin{exercise}[$\star\star\star$]\label{exo:b2:structure-determination:10}
An unknown \ce{C9H10O} shows an IR band at \qty{1690}{cm^{-1}}, carbon lines at 200.8 (C), 137.0 (C),
132.9 (CH), 128.6 (CH), 128.0 (CH), 31.8 (\ce{CH2}) and 8.3 (\ce{CH3}), and proton signals at 7.95 (m,
2H), 7.40--7.55 (m, 3H), 2.98 (q, 2H) and 1.22 (t, 3H) (exercise data). Identify it, and explain the
low carbonyl wavenumber.
\end{exercise}

\begin{exercise}[$\star\star\star$]\label{exo:b2:structure-determination:11}
Four isomers \ce{C5H10O2}: pentanoic acid, methyl butanoate, ethyl propanoate and propyl ethanoate.
Give for each one feature of its IR or proton NMR spectrum that tells it from the other three.
\end{exercise}

\begin{exercise}[$\star\star\star$]\label{exo:b2:structure-determination:12}
A report assigns the line at 166.5~ppm of ethyl benzoate to the ring carbon bearing the ester and
the line at 130.6~ppm to the carbonyl carbon. Show, with the \omterm{def:b2:structure-determination:dept}{DEPT} data and the shift chart, that the
assignment cannot be right.
\end{exercise}

\section{Problem: Four Spectra of Jasmine}

\begin{problem}[{Weekend problem --- the formula from the molecular ion, the infrared carbonyl band,
carbon-13 and DEPT, the proton spectrum, and a structure checked against the fragments}]
\label{pb:b2:structure-determination:1}
A colourless liquid with a flowery odour of jasmine gives the following data. \omterm{def:b2:mass-spec-atomic:spectrum}{Mass spectrum} (NIST
WebBook): $m/z$ 150 (31.7), 151 (3.1), 108 (100), 91 (71.7), 90 (40.6), 43 (37.6). Infrared (exercise
data): strong bands at 1740 and \qty{1230}{cm^{-1}}, weak bands just above \qty{3000}{cm^{-1}}, no
broad band near \qty{3300}{cm^{-1}}. Carbon-13, CDCl$_3$ (PubChem data): 170.70, 136.14, 128.56,
128.24, 66.24, 20.82~ppm. Proton, CDCl$_3$, \qty{90}{MHz}: 7.33 (m, 5H), 5.09 (s, 2H), 2.06 (s, 3H).
% ledger: use:benzylacetate, ms:benzylacetate, c13:benzylacetate, nmr:benzylacetate, ir:CO-ester, iso:C12, iso:C13, mass:C12, mass:H1, mass:O16

\begin{center}
\begin{tikzpicture}
\pgfplotsset{dept/.style={width=0.8\linewidth, height=3.2cm, ycomb, x dir=reverse, ymin=0, ymax=1.15,
  xmin=0, xmax=200, ytick=\empty, axis y line=none, axis x line=bottom, x axis line style={-}, tick label style={font=\tiny},
  title style={font=\scriptsize, at={(0.02,0.82)}, anchor=west}, every axis plot/.append style={thick}}}
\begin{axis}[dept, title={decoupled}, xtick={0,20,40,60,80,100,120,140,160,180,200}]
\addplot[omDef] table[x=shift, y=dec] {figdata/out/bachelor-2-nmr-spectra-c.dat};
\end{axis}
\end{tikzpicture}
\end{center}
% figdata: bachelor-2/nmr-spectra.py part c

\medskip
\noindent\textbf{Part I --- The formula.}
\begin{enumerate}
  \item From the 151/150 ratio, estimate the number of carbons.
  \item What does the \omterm{def:b2:mass-spec-atomic:nitrogen-rule}{nitrogen rule} say?
  \item With nine carbons, what remains of the mass 150? Propose a formula with oxygen.
  \item Why is \ce{C10H14O}, also of nominal mass 150, less likely?
  \item Compute the degree of unsaturation.
  \item Compute the \omterm{def:b2:mass-spec-atomic:exact-mass}{monoisotopic mass} of the formula retained.
\end{enumerate}

\medskip
\noindent\textbf{Part II --- Infrared.}
\begin{enumerate}[resume]
  \item Assign the band at \qty{1740}{cm^{-1}}. Which family does its position suggest?
  \item Assign the band at \qty{1230}{cm^{-1}}.
  \item What do the weak bands above \qty{3000}{cm^{-1}} indicate?
  \item What does the absence of a broad band near \qty{3300}{cm^{-1}} exclude?
  \item Is the carbonyl conjugated with the ring? Explain from its wavenumber.
\end{enumerate}

\medskip
\noindent\textbf{Part III --- Carbon-13 and \omterm{def:b2:structure-determination:dept}{DEPT}.}
\begin{enumerate}[resume]
  \item Assign each line to a kind of carbon.
  \item Which line points down in DEPT-135?
  \item Which lines vanish in DEPT-135?
  \item Which lines remain in DEPT-90?
  \item How many kinds of carbon does a monosubstituted benzene ring have? Explain why only two \omterm{def:b2:huckel:aromatic}{aromatic}
        CH lines appear.
  \item The line at 128.24 is the tallest. Does it stand for the most carbons?
\end{enumerate}

\medskip
\noindent\textbf{Part IV --- Proton NMR and structure.}
\begin{enumerate}[resume]
  \item Assign the three proton signals.
  \item Why are the signals at 5.09 and 2.06 singlets?
  \item Why is the \ce{CH2} at 5.09 so far downfield?
  \item Assemble the structure and name it.
  \item Its isomer methyl 2-phenylethanoate, \ce{C6H5CH2COOCH3}, has the same formula. Which datum rules it
        out?
  \item Explain the fragments at 108, 91 and 43.
  \item How many lines would a perfectly resolved decoupled \ce{^{13}C} spectrum of the compound show, and
        how many of them point down in DEPT-135?
\end{enumerate}
\end{problem}
